% ************************** Thesis Abstract *****************************
% Use `abstract' as an option in the document class to print only the titlepage and the abstract.

\chapter*{\huge \leftline{摘要}}

在复杂网络中预测缺失或未来的连接（即链路预测）是社交网络分析、推荐系统和生物信息学中的一项核心任务。传统的启发式方法完全依赖于局部邻域相似性指标，但无法捕获长距离的拓扑关系；另一方面，以 DeepWalk 为代表的网络表示学习模型通过随机游走学习连续向量嵌入，但通常采用均匀的转移概率，忽略了网络的局部聚类结构。

为了解决这些局限性，本文提出了一种用于链路预测的 Jaccard 加权随机游走图表示框架。我们定义了由邻居间 Jaccard 相似度引导的参数化转移概率。通过引入结构平滑因子，模型可以平衡局部聚类开发和全局拓扑探索。节点嵌入通过 Skip-gram 模型进行训练，并在二分类链路预测任务上进行评估。在标准基准数据集（Cora、Citeseer 和 Facebook）上的实验结果表明，我们提出的相似度加权随机游走框架在 AUC-ROC 和平均精度 (AP) 指标上一致优于原始 DeepWalk 和传统启发式基线。\newline
\\
\textbf{关键词: } \textit{链路预测, 社交网络, 网络表示学习, 节点相似度}

